% Unit 3 review sheet -- 2 pages. Heaviest unit, so the tool list is longer.
\documentclass{apchem-worksheet}

\apcunit{3}
\apcunittitle{Properties of Substances and Mixtures}
\apcdoctitle{Review Sheet}
\apcsubtitle{18--22\% of the AP exam \textbullet\ exam Wed Sep 30 (double block)}

\begin{document}
\apctitleblock

\section*{\sffamily The seven tools}

\renewcommand{\arraystretch}{1.45}
\noindent\begin{tabular}{@{}p{3.8cm}p{9.9cm}@{}}
\toprule
\textbf{IMF identification} & LDF in everything; dipole--dipole if polar;
  H-bonding only if H is on N, O, or F; ion--dipole for salts in polar
  solvent. \emph{Compare size first when molecules differ greatly}\\
\textbf{IMF consequences} & stronger IMF $\to$ higher bp, viscosity, surface
  tension; \emph{lower} vapor pressure\\
\textbf{Solid classification} & ionic / metallic / covalent network /
  molecular --- predict mp, hardness, conductivity from the particle type\\
\textbf{Gas laws} & $PV = nRT$ ($R = 0.08206$), kelvin always;
  $P_A = \chi_A P_{\text{tot}}$\\
\textbf{KMT} & $KE_{\text{avg}} \propto T$ only; $u \propto 1/\sqrt{M}$;
  hotter $=$ broader, flatter, right-shifted, same area\\
\textbf{Solutions} & $M = $ mol solute / L \emph{solution};
  $M_1V_1 = M_2V_2$; ionic solutes drawn fully dissociated in ratio\\
\textbf{Beer--Lambert} & $A = \epsilon bc$; calibration slope $= \epsilon b$;
  unknown $c = A \div \text{slope}$\\
\bottomrule
\end{tabular}

\section*{\sffamily Two deviation stories, kept straight}

\begin{itemize}[leftmargin=*,itemsep=0.2em]
  \item \textbf{$PV/nRT < 1$} --- \emph{attractions} dominate (low $T$,
        moderate $P$). Molecules pull on each other, softening wall
        collisions, so measured $P$ is low.
  \item \textbf{$PV/nRT > 1$} --- \emph{particle volume} dominates (high
        $P$). The gas cannot compress as far as the ideal law predicts.
  \item Most ideal: \textbf{high $T$, low $P$}. Most deviant: strong IMFs
        (\ce{H2O}, \ce{NH3}) at low $T$, high $P$.
\end{itemize}

\section*{\sffamily Separation: name the property, not just the technique}

filtration $\to$ particle size \textbullet\ distillation $\to$ boiling point
\textbullet\ chromatography $\to$ relative attraction to stationary vs
mobile phase

\section*{\sffamily The traps, one line each}

\begin{itemize}[leftmargin=*,itemsep=0.22em]
  \item Ranking by force \emph{type} alone --- \ce{BI3} outboils \ce{NH3}
        because dispersion scales with electron count.
  \item Celsius in a gas law. Convert first, every time.
  \item ``Melting breaks the covalent bonds'' --- melting a molecular solid
        breaks only IMFs.
  \item ``Real gases deviate because they're real'' --- name volume
        \emph{or} attractions, and say which direction.
  \item Drawing ionic solutions as intact units instead of separated ions
        in the correct ratio.
  \item Multiplying by the Beer--Lambert slope instead of dividing.
  \item Forgetting $\nu = c/\lambda$ --- the equations sheet gives only
        $E = h\nu$.
\end{itemize}

\section*{\sffamily Self-check --- 10 questions, exam conditions}

\begin{enumerate}[leftmargin=*,itemsep=0.4em]
  \item All IMFs in \ce{CH3NH2}:
        \answerblank[6.3cm]{LDF, dipole--dipole, H-bonding}
  \item Higher bp, \ce{H2O} or \ce{H2Te}? Why?
        \answerblank[5.5cm]{\ce{H2O} --- hydrogen bonding}
  \item Solid type: conducts as a solid, malleable:
        \answerblank[3.5cm]{metallic}
  \item $V$ of \SI{2.00}{\mole} at \SI{1.00}{\atm}, \SI{273}{\kelvin}:
        \answerblank[3cm]{\SI{44.8}{\liter}}
  \item $P_{\ce{Ar}}$ if $\chi = 0.30$ and $P_{\text{tot}} =
        \SI{5.0}{\atm}$: \answerblank[3cm]{\SI{1.5}{\atm}}
  \item Faster at equal $T$: \ce{CH4} or \ce{SO2}?
        \answerblank[2.5cm]{\ce{CH4}}
  \item $[\ce{NO3-}]$ in \SI{0.20}{\Molar} \ce{Mg(NO3)2}:
        \answerblank[3cm]{\SI{0.40}{\Molar}}
  \item Dilute \SI{15.0}{\milli\liter} of \SI{2.00}{\Molar} to
        \SI{60.0}{\milli\liter}: \answerblank[3cm]{\SI{0.500}{\Molar}}
  \item Technique to separate two miscible liquids:
        \answerblank[3.5cm]{distillation}
  \item $A = 0.240$, slope $= \SI{60.0}{\per\Molar}$. Find $c$:
        \answerblank[3cm]{\SI{0.00400}{\Molar}}
\end{enumerate}

\begin{apcnote}[How to study for this exam]
This is a double block and the largest unit --- budget two study sessions,
not one. Session 1: IMFs, solids, phases (WS 1). Session 2: gases, KMT,
deviations (WS 2), then solutions and spectroscopy (WS 3--5). Redo WS 5's
FRQ cold; the exam's short FRQ is its twin. For every ``explain'' question,
name the particles and the force before naming the property.
\end{apcnote}

\end{document}
